% /Users/pikachu/books/payments-book/src/parts/part4/ch20-subscriptions.tex
\chapter{Подписки и Card-on-File}
\label{ch:subscriptions}

\begin{kaobox}[title=Что вы узнаете из этой главы]
  Подписка~--- это не~просто повторяющееся списание. Это~цепочка согласованных
  транзакций с~собственным протоколом, нарушение которого приводит к~штрафам
  от~схем, отказам авторизации и~потере клиентов. Из~этой главы вы~узнаете,
  как устроен Stored Credential Framework и~почему он~появился; в~чём
  принципиальная разница между CIT и~MIT с~точки зрения авторизационного
  flow; зачем нужен Network Transaction~ID и~что~происходит при~потере
  его~цепочки; как~правильно построить retry при~мягком отказе MIT; и~как
  Account Updater снижает involuntary churn.
\end{kaobox}


\section{Stored Credential Framework}
\label{sec:subscriptions-scf}

До~2017~года мерчанты с~подписочной моделью решали вопрос повторных
списаний каждый по-своему: хранили PAN клиента (или~поручали PSP),
раз~в~месяц отправляли авторизационный запрос и~надеялись на~лучшее.
Никакого стандартизированного способа сообщить эмитенту \enquote{эта
  транзакция проводится по~ранее данному согласию клиента} не~существовало.
Эмитент видел транзакцию, инициированную без~присутствия клиента, и~нередко
отклонял её~как~подозрительную или~применял более строгий антифрод-scoring.

\scheme{Visa} ввела Stored Credential Framework~(SCF) в~октябре 2017~года,
\scheme{Mastercard}~--- в~том~же цикле обновлений~\sidecite{visa-scf}\sidecite{mc-scf}.
Идея фреймворка проста: каждая транзакция с~сохранёнными реквизитами
карты помечается специальными флагами в~авторизационном запросе, которые
сообщают эмитенту контекст~--- является~ли это~первым согласием клиента
или~последующим списанием. Эмитент получает возможность применять
соответствующую авторизационную политику, а~не~единый параноидальный
фильтр для~всех CNP-транзакций.

В~ISO~8583 контекст SCF передаётся через несколько полей: поле
Stored Credential Usage Indicator (используется ли~первичное или~последующее
сохранённое согласие), Stored Credential Type Indicator (рекуррентная
подписка, рассрочка, unscheduled~--- нерегулярное с~переменной суммой,
industry practice~--- отраслевые нормы), и~по~завершении первоначальной
транзакции~--- Network Transaction~ID, который становится идентификатором
всей дальнейшей цепочки~\sidecite{visa-scf}\sidecite{mc-scf}.


\section{CIT vs MIT: различия и~последствия}
\label{sec:subscriptions-cit-mit}

Разграничение между customer-initiated transaction~(CIT) и~merchant-initiated
transaction~(MIT) кажется очевидным на~первый взгляд: первую инициирует
клиент, вторую~--- мерчант без~активного участия клиента. Но~последствия
этого разграничения для~авторизационного flow значительны, а~ошибка
в~классификации~--- дорогостояща.

При CIT~--- customer-initiated transaction~--- клиент присутствует
при~оплате: нажимает кнопку \enquote{оплатить}, сканирует карту, вводит
CVV. К~таким транзакциям применяется SCA (Strong Customer Authentication
по~PSD2 для~европейского трафика), для~них обязателен или~рекомендован 3DS. Именно CIT устанавливает
согласие клиента на~хранение реквизитов и~является точкой входа в~SCF-цепочку:
в~ответе на~успешную CIT эмитент возвращает Network Transaction~ID~(NTID),
который мерчант обязан сохранить.

Зеркальная ситуация~--- MIT, merchant-initiated transaction: клиент
не~присутствует. Это~ежемесячное списание подписки в~04:00, платёж
за~доставку по~сохранённой карте, post-auth корректировка суммы.
К~MIT SCA не~применяется~--- клиент уже дал согласие при~CIT. Однако
MIT обязан содержать в~запросе
NTID первоначальной CIT: без~него эмитент не~распознаёт транзакцию
как~часть согласованной цепочки и~вправе применить более строгую политику
отказа~\sidecite{visa-scf}.

\todofig{Decision tree «CIT или MIT?». Корень: «Инициирована ли транзакция
  активным действием клиента прямо сейчас?». Ветвь «Да» → CIT. Под CIT:
  «3DS применим?» → «Да» → «3DS обязателен / рекомендован». «Нет» → CVV/AVS.
  Результат CIT: «сохранить NTID». Ветвь «Нет» → MIT. Под MIT:
  «Тип» → 3 ветви: Recurring (регулярная, фиксированная сумма), Installment
  (рассрочка), Unscheduled (нерегулярная, переменная). Требование для всех
  MIT: «передать NTID из CIT + флаги MIT». Справа: «Ошибка: MIT без NTID
  → эмитент может отклонить или применить SCA-требование». Стиль:
  минималистичный, с цветовой кодировкой CIT (синий) / MIT (оранжевый).}
% \label{fig:subscriptions-cit-mit}

Ошибка классификации~--- передать MIT-транзакцию без~соответствующих
флагов или~с~флагами CIT~--- приводит к~двум видам проблем. Первый~---
регуляторный: в~Европе MIT без~корректного SCA-exemption нарушает PSD2,
что~влечёт ответственность эмитента и~потенциальные штрафы. Второй~---
операционный: эмитент, не~распознавший MIT, применяет антифрод-политику
для~незнакомой CNP-транзакции и~отклоняет её~с~кодом \rc{05} там, где
мог~бы одобрить.


\section{Network Transaction~ID и~chaining}
\label{sec:subscriptions-ntid}

Network Transaction~ID~(NTID)~--- это~уникальный идентификатор, который
схема присваивает каждой успешно авторизованной транзакции и~возвращает
в~авторизационном ответе~\sidecite{visa-scf}. Для~SCF-цепочки NTID первой
CIT~--- это~якорь, к~которому привязываются все последующие MIT. Каждая
следующая MIT-транзакция передаёт NTID предыдущей, создавая цепочку,
по~которой эмитент может проследить историю согласия клиента.

Потеря NTID~--- частая причина деградации authorization rate у~мерчантов
при~смене PSP или~при~миграции процессинговой инфраструктуры. Сценарий
выглядит так: мерчант переходит к~другому PSP, переносит сохранённые
реквизиты карт, но~не~переносит NTID первоначальных CIT-транзакций.
Следующий биллинговый цикл проходит без~NTID~--- и~часть эмитентов,
особенно европейских с~жёсткой PSD2-политикой, начинает отклонять
MIT-транзакции, требуя либо SCA, либо NTID. Сохранение NTID при~миграции~---
обязательный пункт технического задания при~любом переходе между
процессорами.

Дополнительная ценность NTID-цепочки открылась с~появлением Compelling
Evidence~3.0: как~мы~обсудили в~главе~\ref{ch:disputes}, история транзакций
с~тем~же устройством и~NTID-цепочкой служит доказательством при~fraud-диспуте
(код~10.4). Billing engine, который правильно хранит и~передаёт NTID,~---
это одновременно и~система снижения churn, и~система защиты от~friendly
fraud.

\begin{practice}
  Храните NTID в~своей системе отдельно от~токена или~PAN карты:
  это~разные идентификаторы с~разными жизненными циклами. NTID привязан
  к~конкретному согласию клиента; токен~--- к~реквизитам карты. При~перевыпуске
  карты токен обновится автоматически (или~через Account Updater), но~NTID
  останется тем~же~--- он~по-прежнему идентифицирует исходное согласие.
  При~миграции к~другому PSP убедитесь, что~в~техническом задании явно
  прописана передача NTID для~каждой активной подписки.
\end{practice}


\section{Soft decline при~MIT: обязательный retry protocol}
\label{sec:subscriptions-soft-decline}

Когда MIT-транзакция получает soft decline с~кодом \rc{1A}~--- SCA
Required,~--- это~специальный сигнал, появившийся после введения PSD2:
эмитент требует прохождения 3DS для~данной транзакции~\sidecite{visa-scf}.
Для~мерчанта вне~зоны PSD2 этот код встречается редко; для~европейского
трафика~--- с~определённой периодичностью, особенно у~клиентов, давно
не~авторизовавшихся в~банке.

Правильная реакция на~\rc{1A}: не~пытаться провести транзакцию повторно
без~изменений (это~запрещено правилами схем и~бессмысленно), а~отправить
клиенту ссылку на~3DS-аутентификацию в~удобный момент. После успешного
прохождения 3DS клиентом мерчант получает новый NTID и~может возобновить
MIT-цепочку~\sidecite{visa-scf}.

Для~остальных soft decline при~MIT~--- \rc{05}, \rc{51}~--- действует
схемный retry protocol, который отличается от~произвольного retry
для~разовых транзакций~\sidecite{visa-retry-rules}. Суть схемного
протокола~--- строго ограниченное число попыток в~заданные тайм-слоты
с~обязательным указанием флага \enquote{это~retry}. Повторная попытка
без~флага retry классифицируется схемой как~новая MIT без~NTID~--- нарушение,
которое может привести к~штрафам.

\begin{important}
  Произвольный retry при~MIT вне~схемного протокола~--- прямое нарушение
  правил Visa и~Mastercard. Это~не~просто этическая норма: нарушение
  обнаруживается автоматически на~уровне авторизационного процессинга
  схемы и~влечёт штрафы эквайеру, который они~перекладывают на~мерчанта.
  Убедитесь, что~ваш PSP реализует retry по~схемному протоколу, а~не~самодельной
  логике.
\end{important}

Подробнее об~общей retry logic и~классификации response codes~---
в~главе~\ref{ch:decline-management}.


\section{Account Updater: VAU и~ABU}
\label{sec:subscriptions-account-updater}

Карты перевыпускаются. Эмитент выдаёт новую карту с~новым PAN при~компрометации,
с~новой датой истечения при~плановом перевыпуске, с~новыми реквизитами
при~смене продуктовой программы. Мерчант с~сохранёнными реквизитами
об~этом не~знает и~продолжает пытаться провести MIT по~старому PAN~---
получая \rc{54} Expired~Card или~\rc{14} Invalid~Account Number
и~теряя клиента.

Account Updater решает эту проблему на~уровне схемы. Visa предоставляет
сервис Visa Account Updater~(VAU), Mastercard~---
Automatic Billing Updater~(ABU)~\sidecite{visa-core-rules}\sidecite{mc-rules}.
Механизм работает через batch-запросы: мерчант периодически~--- обычно
раз~в~неделю или~накануне биллингового цикла~--- передаёт список токенов
или~PAN, для~которых хочет проверить актуальность; схема возвращает
обновлённые реквизиты для~тех~карт, которые были перевыпущены и~для~которых
эмитент предоставил обновление.

Несколько важных ограничений. Первое: эмитент участвует добровольно;
если эмитент не~подключён к~VAU/ABU, данные не~обновятся. В~России
и~СНГ покрытие эмитентов Account Updater значительно ниже, чем~в~США.
Второе: обновление происходит не~мгновенно, а~batch-цикл имеет задержку~---
для~карт, перевыпущенных за~день до~биллинга, обновление может не~успеть.
Третье: если мерчант использует сетевые токены~(VTS/MDES), Account Updater
по~большей части не~нужен~--- токен обновляется схемой автоматически
и~прозрачно для~мерчанта. Именно поэтому переход на~network tokens~---
более надёжное долгосрочное решение проблемы перевыпуска карт,
чем~batch Account Updater~\sidecite{visa-vts}\sidecite{mc-mdes}.


\section{Процесс подписки без~технического churn}
\label{sec:subscriptions-anti-churn}

\textit{Technical churn}~--- это~потеря клиента-подписчика не~из-за~его
желания отменить сервис, а~из-за~технических сбоев в~биллинге: decline
по~перевыпущенной карте, неправильный retry, потеря NTID, пропуск
SCA-требования. В~зрелых подписочных бизнесах technical churn составляет
значительную часть общего оттока, нередко сопоставимую с~voluntary churn.
Устранение технического churn~--- в~значительной мере инженерная задача,
поддающаяся решению.

Правильный процесс подписки выглядит следующим образом. На~этапе онбординга~---
первичная CIT-транзакция с~3DS-аутентификацией: это~одновременно
и~авторизация первого платежа, и~установление согласия клиента в~рамках
SCF. Из~ответа на~CIT обязательно сохраняются NTID и~токен (или~PAN,
если токенизация не~используется). Клиент получает подтверждение,
в~котором явно описаны условия подписки~--- дата, сумма, периодичность.

На~каждом биллинговом цикле MIT-транзакция отправляется с~корректными
флагами (Stored Credential Type~= Recurring, NTID из~первой CIT, флаг
MIT). Накануне биллинга~--- запрос к~VAU/ABU для~обновления реквизитов;
если используются сетевые токены~--- этот шаг избыточен. При~получении
мягкого decline~--- retry по~схемному протоколу, но~не~произвольный.
При~\rc{1A}~--- уведомление клиенту с~запросом на~3DS. При~жёстком
decline~--- уведомление клиенту о~необходимости обновить реквизиты.

Отдельная составляющая~--- dunning, последовательность коммуникаций
с~клиентом при~неудачном биллинге. Исследования показывают, что~оптимальная
точка первого уведомления~--- не~в~момент decline, а~за~2--3~дня до~следующей
попытки: клиент, получивший напоминание о~необходимости пополнить счёт
или~обновить карту, с~большей вероятностью сделает это~до~retry~\todocite{данные по~эффективности dunning-стратегий в~подписочных сервисах}.

\begin{practice}
  Ключевые метрики технического churn: involuntary churn rate
  (доля отписок, вызванных техническими отказами биллинга), authorization
  rate для~MIT отдельно от~CIT, доля \rc{54} в~отказах (индикатор неработающего
  Account Updater или~отсутствия network tokens), доля \rc{1A} (индикатор
  проблемы с~NTID или~PSD2-exemption). Мониторинг этих метрик в~разрезе
  когорт по~дате онбординга позволяет отделить системную проблему
  от~случайного всплеска.
\end{practice}


\section*{Итоги главы}

\begin{itemize}
  \item Stored Credential Framework~(SCF) стандартизировал передачу
        контекста сохранённых реквизитов: эмитент знает, является~ли
        транзакция первичным согласием клиента или~последующим списанием
        по~нему.
  \item CIT~(customer-initiated) устанавливает согласие и~порождает
        NTID; MIT~(merchant-initiated) обязан передавать NTID из~первой
        CIT~--- без~него эмитент применяет более строгую политику отказа.
  \item Потеря NTID при~миграции между PSP~--- распространённая причина
        деградации authorization rate в~подписках; NTID и~токен~---
        разные идентификаторы с~разными жизненными циклами, хранить
        нужно оба.
  \item Soft decline \rc{1A} при~MIT требует 3DS-аутентификации клиента,
        а~не~повторной попытки транзакции; произвольный retry вне~схемного
        протокола запрещён правилами схем и~влечёт штрафы.
  \item Account Updater~(VAU/ABU) снижает \rc{54} Expired~Card через
        batch-обновление реквизитов; сетевые токены решают ту~же проблему
        автоматически и~без~batch-запросов.
  \item Technical churn устраняется инженерными методами: правильный
        SCF-flow, NTID-chaining, network tokens, схемный retry
        protocol и~своевременный dunning~--- в~совокупности дают
        2--4~п.п. снижения involuntary churn.
\end{itemize}
